\subsection{An exact comparison for local decoupling terminal costs}
\label{decprofile:statement}

This section concerns an abstract profile optimization. Its proof is purely
geometric and does not use the preceding analytic estimate.

For any 1-Lipschitz \(g:[0,1]\to\mathbb R\), define
\[
 c_g(m,n)=g(m)-\min_{[m,n]}g.
\]
An edge \(n\to m<n\) is admissible when \(2n-m\le1\).
Let \(\Phi_g(n)\) be the infimum of the costs of finite admissible chains
from \(n<1\) to zero, and write
\(\Phi_g(1)=\lim_{n\uparrow1}\Phi_g(n)\). The limit and the finite-chain
properties used below are justified in the proof.
For \(0\le m<1\), set
\begin{equation}\label{decprofile:terminal}
 L=1-m,\qquad n=\frac{1+m}{2},\qquad
 T_6(g,m)=\frac{2c_g(m,n)+L-g(1)+g(m)}3.
\end{equation}

\begin{theorem}[The terminal expression is attained by ordinary chains]
\label{decprofile:identity}
For every 1-Lipschitz \(g:[0,1]\to\mathbb R\),
\begin{equation}\label{decprofile:exact}
 \begin{aligned}
 \Phi_g(1)&\le\Phi_g(m)+T_6(g,m)\qquad(0\le m<1),\\
 \inf_{0\le m<1}\bigl(\Phi_g(m)+T_6(g,m)\bigr)&=\Phi_g(1).
 \end{aligned}
\end{equation}
No sign condition, dimension barriers, or endpoint normalization are
required. At every fixed start \(n_0<1\), an actual chain with
\(O(1+\log(1/(1-n_0)))\) edges has cost at most the displayed upper
bound. Uniform profile
perturbations of size \(\varepsilon\) and grid rounding of mesh \(\rho\)
give error \(O(K(\varepsilon+\rho))\), where \(K\) is that edge bound.
\end{theorem}

The theorem rules out an improvement from replacing the terminal portion
of an optimized chain by exactly \eqref{decprofile:terminal}. It does not
rule out other uses of local decoupling.

\subsection{A total-drop potential}
\label{decprofile:potential}

For a 1-Lipschitz function \(f:[0,1]\to\mathbb R\) with \(f(0)=0\),
let \(\mathbf T(f)\) be the infimum of
\[
 \sum_{j\ge1}\left(f(t_j)-\min_{[t_j,t_{j-1}]}f\right)
\]
over good partitions \(1=t_0>t_1>\cdots\to0\), with
\(t_{j-1}\le2t_j\). We claim that
\begin{equation}\label{decprofile:totaldrop}
 \mathbf T(f)\le\frac{1-2\beta}{3},
 \qquad \beta=\min_{[1/2,1]}f.
\end{equation}

First suppose that \(f\) is piecewise affine with finitely many pieces.
Its hard-point set
\[
 H=\left\{t\in[0,1]: f(t)=\min_{[t/2,t]}f\right\}
\]
is a finite union of closed interval components, possibly degenerate.
It contains zero, and \(f\) is nonincreasing on each component.
Keleti--Shmerkin's identity \cite[Lemma 5.4]{ks} gives
\begin{equation}\label{decprofile:hardidentity}
 \mathbf T(f)=
 \sum_{[u,v]\text{ component of }H}\bigl(f(u)-f(v)\bigr).
\end{equation}
Consider the potential
\[
 P(t)=\frac{t-2f(t)}3.
\]
On a hard component \([u,v]\), its drop
\(\delta=f(u)-f(v)\) satisfies \(0\le\delta\le v-u\). Hence
\begin{equation}\label{decprofile:component}
 P(v)-P(u)=\frac{v-u+2\delta}{3}\ge\delta.
\end{equation}
Now take the endpoints \(v<u\) of a gap between components, and put
\(\Delta=f(u)-f(v)\). If \(\Delta\le0\), then \(P(u)\ge P(v)\).
If \(\Delta>0\), hardness at \(u\) forces \(v<u/2\), since otherwise
\(v\in[u/2,u]\) would give \(f(v)\ge f(u)\). Also
\(f(u/2)\ge f(u)\), and therefore
\[
 0<\Delta\le f(u/2)-f(v)\le u/2-v.
\]
It follows that
\begin{equation}\label{decprofile:gap}
 P(u)-P(v)=\frac{u-v-2\Delta}{3}\ge v/3\ge0.
\end{equation}
Let \(r=\max H\). Summing
\eqref{decprofile:component} and \eqref{decprofile:gap}, and using
\(P(0)=0\), proves
\[
 \mathbf T(f)\le P(r)=\frac{r-2f(r)}3.
\]
If \(r\ge1/2\), then \(f(r)\ge\beta\) and
\(r-2f(r)\le1-2\beta\). If \(r<1/2\), Lipschitz continuity gives
\[
 f(r)\ge f(1/2)-(1/2-r)\ge\beta-1/2+r,
 \qquad r-2f(r)\le1-2\beta-r\le1-2\beta.
\]
This proves \eqref{decprofile:totaldrop} in the piecewise affine case.
The finite-chain approximation below extends it to every 1-Lipschitz
function, without an assumed continuity theorem for infinite partitions.

\subsection{Reversal, concatenation, and finite approximation}
\label{decprofile:proof}

Let \(h:[0,1]\to\mathbb R\) be 1-Lipschitz with \(h(0)=0\),
and put \(e=h(1)\) and \(f(t)=h(1-t)-e\).
Under reversal, an original edge \(y\to x\) corresponds to a ratio-two
interval \([1-y,1-x]\). The difference between its original cost and
the total-drop summand is the change in \(f\) between the two endpoints.
Telescoping, and then letting the start approach one, yields
\begin{equation}\label{decprofile:reversal}
 \Phi_h(1)=\mathbf T(f)-e.
\end{equation}
For completeness, a finite ratio-two partition starting at \(\zeta>0\)
can be extended towards zero by successive halvings with additional cost
at most \(\zeta\). Conversely, an infinite good partition can be truncated
at its interval crossing \(\zeta\), moving that endpoint to \(\zeta\);
the crossing cost changes by at most \(2\zeta\). These observations
justify the limit in \eqref{decprofile:reversal}.
Since
\[
 \min_{[1/2,1]}f=\min_{[0,1/2]}h-e,
\]
equations \eqref{decprofile:totaldrop} and \eqref{decprofile:reversal}
give
\begin{equation}\label{decprofile:local}
 \Phi_h(1)\le\frac{1-e-2\min_{[0,1/2]}h}{3}.
\end{equation}

We record the finite-chain details before applying this bound.
Truncating a chain at its first crossing of a lower start can only
decrease its cost. Conversely, successive maximal admissible edges bridge
two starts below one with total cost at most their distance. Thus the
finite-start infima are nondecreasing, are bounded by one, and satisfy
\begin{equation}\label{decprofile:start}
 0\le\Phi_h(1)-\Phi_h(1-\zeta)\le\zeta.
\end{equation}
For \(l<m<n\), comparison of the two interval minima gives
\begin{equation}\label{decprofile:merge}
 c_h(l,n)\le c_h(l,m)+c_h(m,n).
\end{equation}
Merge consecutive edges whenever their union is admissible.
In a resulting chain no two consecutive edges can be merged, so the
complementary distance to one more than doubles every two edges.
Its number of edges is at most
\(2\lceil\log_2(1/\zeta)\rceil+O(1)\).
An infimizing sequence of these bounded endpoint lists has a convergent
subsequence. Admissibility is closed, and interval minima depend
continuously on endpoints. Removing repeated limiting endpoints gives
a chain attaining the finite-start infimum.

For a general \(h\), approximate it uniformly by its piecewise affine
interpolants, with the same endpoint values and Lipschitz constant.
At any fixed finite start, the piecewise affine bound supplies chains
of uniformly bounded length. Pass to a convergent subsequence of their
endpoint lists. The interval minima and costs converge and admissibility
is preserved. This proves \eqref{decprofile:local} for every finite start;
\eqref{decprofile:start} gives its limiting form. It also establishes
\eqref{decprofile:totaldrop} in general via the reversal identity.

Apply \eqref{decprofile:local} to
\[
 h(x)=\frac{g(m+Lx)-g(m)}{L},\qquad L=1-m.
\]
Admissibility rescales exactly:
\[
 2(m+Ly)-(m+Lx)\le1\quad\Longleftrightarrow\quad 2y-x\le1.
\]
Multiplication by \(L\) gives the upper cost \(T_6(g,m)\) for a tail
chain ending at \(m\). Concatenate it with a chain from \(m\) to zero.
This proves the first assertion of \eqref{decprofile:exact}.
The Lipschitz bound gives
\[
 0\le T_6(g,m)\le1-m.
\]
Letting \(m\uparrow1\) proves the opposite inequality for the infimum,
and hence the exact identity.

For a finite start \(n_0>m\), the rescaled tail has
\[
 O\left(1+\log\frac{1-m}{1-n_0}\right)
\]
edges. Compress the concatenated chain once more to obtain the edge
bound in Theorem~\ref{decprofile:identity}. Starts \(n_0\le m\) are
handled by monotonicity. A uniform perturbation of \(g\) by
\(\varepsilon\) changes each fixed edge cost by at most \(2\varepsilon\).
Rounding endpoints down to a grid of mesh \(\rho\), with zero, one and
the prescribed finite start on the grid, preserves admissibility and changes each edge cost by at
most \(2\rho\). Indeed
\[
 \left\lfloor m/\rho\right\rfloor
 \ge\left\lfloor(2n-1)/\rho\right\rfloor
 \ge2\left\lfloor n/\rho\right\rfloor-1/\rho
\]
when \(1/\rho\) is integral. This proves the stated stability.

\subsection{The weighted refined-decoupling family}
\label{decprofile:weighted}

The conditional terminal expression associated with exponents
\(2\le p\le6\) is
\[
 T_p(g,m)=\frac2p I+
       \left(\frac2p-\frac13\right)W+
       \left(1-\frac2p\right)c,
\]
where, with \(L,n\) as in \eqref{decprofile:terminal},
\[
 I=L+g(m)-g(1),\qquad
 W=-L/2+g(1)-g(n),\qquad
 c=g(m)-\min_{[m,n]}g.
\]
At \(p=6\) this equals \eqref{decprofile:terminal}. Direct algebra gives
\begin{align*}
 T_p(g,m)-T_6(g,m)
 &=\left(\frac2p-\frac13\right)(I+W-c)\\
 &=\left(\frac2p-\frac13\right)
       \left(\frac L2-g(n)+\min_{[m,n]}g\right)\ge0.
\end{align*}
The last inequality is exactly the Lipschitz bound over an interval of
length \(L/2\). Consequently
\[
 \inf_{\substack{0\le m<1\\2\le p\le6}}
       \bigl(\Phi_g(m)+T_p(g,m)\bigr)=\Phi_g(1).
\]
Thus allowing this entire weighted exponent family leaves the abstract
profile optimum unchanged.
